\begin{example}  % 动点引入，距离公式
% 入题思路例题1.1
（2026海淀一模）已知椭圆 $E:\frac{x^2}{a^2}+\frac{y^2}{b^2}=1\ (a>b>0)$ 的离心率为 $\frac{\sqrt{3}}{2}$，其左、右焦点和上顶点构成的三角形面积为 $3\sqrt{3}$.\\
（1）求椭圆 $E$ 的方程；\\
（2）设 $O$ 为坐标原点，$M(x_0,y_0)$ 为第一象限内 $E$ 上的动点，点 $N$ 在直线 $y=-2$ 上，且 $OM\perp ON$. 过 $A(0,3)$ 作 $MN$ 的垂线交直线 $OM$ 于 $B$，求 $\lvert AB\rvert$ 的值.
\end{example}

\begin{example} % 直线引入
% 入题思路例题1.6
（2026丰台一模）已知椭圆 $E:\frac{x^2}{a^2}+\frac{y^2}{b^2}=1\ (a>b>0)$，椭圆 $E$ 上一点 $A(2,1)$ 到两焦点的距离之和为 $2\sqrt{6}$.\\
（1）求椭圆 $E$ 的方程和离心率；\\
（2）过点 $P(2,3)$ 的直线 $l$ 与椭圆 $E$ 交于不同的两点 $B,C$. 若直线 $AB$ 与直线 $AC$ 的斜率之和为 $0$，求直线 $l$ 的方程.
\end{example}

\begin{example} % 直线引入，直线相交求交点，三点共线
    % 斜率例题
（2026朝阳一模）已知椭圆 $E:\frac{x^2}{a^2}+\frac{y^2}{b^2}=1\ (a>b>0)$ 的离心率为 $\frac{\sqrt{3}}{2}$，$A,B$ 分别为椭圆 $E$ 的上、下顶点，且 $B(0,-1)$.\\
（1）求椭圆 $E$ 的方程；\\
（2）过点 $D(-1,-1)$ 的斜率存在且不为 $1$ 的直线 $l$ 与椭圆 $E$ 交于不同的两点 $M,N$（均不与点 $A$ 重合），点 $P$ 与点 $M$ 关于原点 $O$ 对称，直线 $BP$ 与直线 $OD$ 交于点 $Q$. 求证：直线 $NQ$ 经过点 $A$.
\end{example}

\begin{example} % 烂题
（2026西城一模）已知椭圆 $E:\frac{x^2}{a^2}+\frac{y^2}{b^2}=1\ (a>b>0)$ 的右顶点为 $A$，下顶点为 $B$，左、右焦点分别为 $F_1,F_2$，且 $\angle F_1BF_2=\frac{2\pi}{3}$，圆 $W:x^2+y^2=1$ 经过点 $B$. $O$ 为坐标原点.\\
（1）求椭圆 $E$ 的方程；\\
（2）设圆 $W$ 的切线 $l$ 与椭圆 $E$ 相交于 $C,D$ 两点，点 $Q(x_Q,y_Q)\ (y_Q>0)$ 在直线 $l$ 上，且满足 $OQ\perp OD$. 过点 $Q$ 作 $QP\perp y$ 轴于点 $P$，求证：$AP\parallel BF_1$.
\end{example}
\begin{example} % 直线，平行四边形
    % 平行四边形例题
（2026东城一模）已知椭圆 $C:\frac{x^2}{a^2}+\frac{y^2}{b^2}=1\ (a>b>0)$ 经过点 $(0,1)$，离心率为 $\frac{\sqrt{6}}{3}$.\\
（1）求椭圆 $C$ 的方程；\\
（2）过点 $D(1,0)$ 作斜率为 $k$ 的直线 $l$ 与椭圆 $C$ 交于 $A,B$ 两点，点 $E(3,n)$ 不在直线 $l$ 上，直线 $AE,BE$ 分别交直线 $x=2$ 于点 $M,N$. 求证：四边形 $DMEN$ 为平行四边形.
\end{example}

\begin{example} % 动点引入（在直线上），直线相交求交点
% 入题思路例题1.2
（2025海淀一模）已知椭圆 $W:\frac{x^2}{a^2}+\frac{y^2}{b^2}=1\ (a>b>0)$，$A,B$ 分别是 $W$ 的左、右顶点，$C$ 是 $W$ 的上顶点，$\triangle ABC$ 的面积为 $2$，且 $\lvert AC\rvert=\sqrt{5}$.\\
（1）求椭圆 $W$ 的方程及长轴长；\\
（2）已知点 $M(2,1)$，点 $P$ 在直线 $AC$ 上，设直线 $PM$ 与 $x$ 轴交于点 $E$，直线 $BP$ 与直线 $EC$ 交于点 $Q$，判断点 $Q$ 是否在椭圆 $W$ 上，并说明理由.
\end{example}

\begin{example} % 动点引入，线段为直径的圆过某点
    % 圆例题
（2025东城一模）已知椭圆 $E:\frac{x^2}{a^2}+\frac{y^2}{b^2}=1\ (a>b>0)$ 过点 $(0,1)$，离心率为 $\frac{\sqrt{6}}{3}$，$E$ 上的点 $A(m,n)\ (n\neq0)$ 关于 $x$ 轴的对称点为 $B$. 设 $O$ 为原点，$\overrightarrow{OH}=\lambda\overrightarrow{OA}\ (0<\lambda<1)$，过点 $H$ 与 $x$ 轴平行的直线交 $E$ 于点 $P,Q$.\\
（1）求椭圆 $E$ 的方程；\\
（2）若点 $B$ 在以 $PQ$ 为直径的圆上，求 $\lambda$ 的值.
\end{example}

\begin{example} % 动点引入，韦达定理求坐标，面积表示，计算量较大
（2025西城一模）已知椭圆 $E:\frac{x^2}{a^2}+\frac{y^2}{b^2}=1\ (a>b>0)$ 的离心率为 $\frac{\sqrt{6}}{3}$，$A$ 为椭圆 $E$ 上一点，且点 $A$ 到椭圆 $E$ 的两个焦点的距离之和等于 $2\sqrt{6}$.\\
（1）求椭圆 $E$ 的方程；\\
（2）若 $A$ 关于原点 $O$ 的对称点为 $B$，过点 $A$ 与 $AB$ 垂直的直线与椭圆 $E$ 的另一个交点为 $C$，$AH\perp x$ 轴于点 $H$，直线 $BC$ 与 $x$ 轴交于点 $M$. 用 $S_{\triangle BOM}$ 与 $S_{\triangle AOH}$ 分别表示 $\triangle BOM$ 与 $\triangle AOH$ 的面积，证明：$S_{\triangle BOM}=2S_{\triangle AOH}$.
\end{example}

\begin{example} % 直线引入，定值定点，非对称韦达
    % 非对称 定点例题
（2025丰台一模）已知椭圆 $E:\frac{x^2}{a^2}+\frac{y^2}{b^2}=1\ (a>b>0)$，以 $E$ 的两个焦点与短轴的一个端点为顶点的三角形是等腰直角三角形，且面积为 $1$.\\
（1）求椭圆 $E$ 的方程；\\
（2）过点 $P(2,0)$ 的直线与椭圆 $E$ 交于不同的两点 $M,N$. 过点 $M$ 作直线 $x=1$ 的垂线，垂足为 $Q$. 求证：直线 $NQ$ 过定点.
\end{example}

\begin{example} % 直线引入，定值定点，非对称韦达
    % 非对称 定值例题
（2025朝阳一模）已知椭圆 $E:\frac{x^2}{a^2}+\frac{y^2}{b^2}=1\ (a>b>0)$ 的右焦点为 $F(1,0)$，离心率为 $\frac{1}{2}$.\\
（1）求椭圆 $E$ 的方程；\\
（2）过点 $M(4,0)$ 作直线 $l$ 与椭圆 $E$ 交于不同的两点 $A,B$. 设 $C\left(\frac{5}{2},\frac{3}{2}\right)$，直线 $BC$ 与直线 $x=1$ 交于点 $N$，求证：直线 $AN$ 的斜率为定值.
\end{example}

\begin{example} % 直线引入，定值定点，长度转圆
    % 练习题
（2024丰台一模）已知椭圆 $E:\frac{x^2}{a^2}+\frac{y^2}{b^2}=1\ (a>b>0)$ 的焦距为 $4\sqrt{2}$，以椭圆 $E$ 的四个顶点为顶点的四边形的周长为 $16$.\\
（1）求椭圆 $E$ 的标准方程；\\
（2）过点 $S(0,1)$ 的直线 $l$ 交椭圆 $E$ 于 $P,Q$ 两点，线段 $PQ$ 的中点为 $M$. 是否存在定点 $D$，使得 $\frac{\lvert DM\rvert}{\lvert PQ\rvert}=\frac{1}{2}$？若存在，求出 $D$ 的坐标；若不存在，请说明理由.
\end{example}

\begin{example} % 直线引入，直线相交求交点，距离转坐标
（2024西城一模）已知椭圆 $G:\frac{x^2}{a^2}+\frac{y^2}{b^2}=1\ (a>b>0)$ 的一个顶点为 $A(-2,0)$，离心率为 $\frac{1}{2}$.\\
（1）求椭圆 $G$ 的方程；\\
（2）设 $O$ 为原点. 直线 $l$ 与椭圆 $G$ 交于 $C,D$ 两点（$C,D$ 不是椭圆的顶点），$l$ 与直线 $x=2$ 交于点 $E$，直线 $AC,AD$ 分别与直线 $OE$ 交于点 $M,N$. 求证：$\lvert OM\rvert=\lvert ON\rvert$.
\end{example}

\begin{example} % 直线引入，圆的切线，平行四边形及面积
    % 练习题
（2024东城一模）已知椭圆 $C:\frac{x^2}{a^2}+\frac{y^2}{b^2}=1\ (a>b>0)$ 的短轴长为 $2\sqrt{3}$，离心率 $e=\frac{\sqrt{2}}{2}$.\\
（1）求椭圆 $C$ 的方程；\\
（2）设 $O$ 为坐标原点，直线 $l$ 是圆 $x^2+y^2=1$ 的一条切线，且直线 $l$ 与椭圆 $C$ 交于 $M,N$ 两点. 若平行四边形 $OMPN$ 的顶点 $P$ 恰好在椭圆 $C$ 上，求平行四边形 $OMPN$ 的面积.
\end{example}

\begin{example} % 烂题
（2024朝阳一模）已知椭圆 $E:\frac{x^2}{a^2}+\frac{y^2}{b^2}=1\ (a>b>0)$ 的离心率为 $\frac{\sqrt{2}}{2}$，$A,B$ 分别是 $E$ 的左、右顶点，$P$ 是 $E$ 上异于 $A,B$ 的点，$\triangle APB$ 的面积的最大值为 $2\sqrt{2}$.\\
（1）求 $E$ 的方程；\\
（2）设 $O$ 为原点，点 $N$ 在直线 $x=2$ 上，$N,P$ 分别在 $x$ 轴的两侧，且 $\triangle APB$ 与 $\triangle NBP$ 的面积相等.\\
（i）求证：直线 $ON$ 与直线 $AP$ 的斜率之积为定值；\\
（ii）是否存在点 $P$ 使得 $\triangle APB\cong\triangle NBP$？若存在，求出点 $P$ 的坐标；若不存在，说明理由.
\end{example}

\begin{example} % 动点引入，距离公式
% 入题思路例题1.3
% 距离例题
（2024海淀一模）已知椭圆 $G:x^2+my^2=m$ 的离心率为 $\frac{\sqrt{2}}{2}$，$A_1,A_2$ 分别是 $G$ 的左、右顶点，$F$ 是 $G$ 的右焦点.\\
（1）求 $m$ 的值及点 $F$ 的坐标；\\
（2）设 $P$ 是椭圆 $G$ 上异于顶点的动点，点 $Q$ 在直线 $x=2$ 上，且 $PF\perp FQ$，直线 $PQ$ 与 $x$ 轴交于点 $M$. 比较 $\lvert MP\rvert^2$ 与 $\lvert MA_1\rvert\cdot\lvert MA_2\rvert$ 的大小.
\end{example}

\begin{example} % 直线，三角形面积
    % 练习题目
（2023丰台一模）已知椭圆 $E:\frac{x^2}{a^2}+\frac{y^2}{b^2}=1\ (a>b>0)$ 的一个顶点为 $A(0,1)$，焦距为 $2$.\\
（1）求椭圆 $E$ 的方程；\\
（2）过点 $P(2,0)$ 的直线与椭圆 $E$ 交于 $B,C$ 两点，过点 $B,C$ 分别作直线 $l':x=t$ 的垂线（点 $B,C$ 在直线 $l'$ 的两侧），垂足分别为 $M,N$. 记 $\triangle BMP,\triangle MNP,\triangle CNP$ 的面积分别为 $S_1,S_2,S_3$，试问：是否存在常数 $t$，使得 $S_1,\frac{1}{2}S_2,S_3$ 总成等比数列？若存在，求出 $t$ 的值；若不存在，请说明理由.
\end{example}

\begin{example} % 直线，
（2023西城一模）已知椭圆 $C:x^2+2y^2=2$，点 $A,B$ 在椭圆 $C$ 上，且 $OA\perp OB$（$O$ 为原点）. 设 $AB$ 的中点为 $M$，射线 $OM$ 交椭圆 $C$ 于点 $N$.\\
（1）当直线 $AB$ 与 $x$ 轴垂直时，求直线 $AB$ 的方程；\\
（2）求 $\frac{\lvert ON\rvert}{\lvert OM\rvert}$ 的取值范围.
\end{example}

\begin{example} % 直线，长度转换
（2023东城一模）已知椭圆 $E:\frac{x^2}{a^2}+\frac{y^2}{b^2}=1\ (a>b>0)$ 的一个顶点为 $A(0,1)$，离心率 $e=\frac{\sqrt{6}}{3}$.\\
（1）求椭圆 $E$ 的方程；\\
（2）过点 $P(-\sqrt{3},1)$ 作斜率为 $k$ 的直线与椭圆 $E$ 交于不同的两点 $B,C$，直线 $AB,AC$ 分别与 $x$ 轴交于点 $M,N$. 设椭圆的左顶点为 $D$，求 $\frac{\lvert MD\rvert}{\lvert MN\rvert}$ 的值.
\end{example}

\begin{example} % 直线，定值定点与非对称韦达
    % 非对称定点例题
（2023朝阳一模）已知椭圆 $E:\frac{x^2}{4}+\frac{y^2}{n}=1\ (0<n<4)$ 经过点 $(\sqrt{2},1)$.\\
（1）求椭圆 $E$ 的方程及离心率；\\
（2）设椭圆 $E$ 的左顶点为 $A$，直线 $l:x=my+1$ 与 $E$ 相交于 $M,N$ 两点，直线 $AM$ 与直线 $x=4$ 相交于点 $Q$. 问：直线 $NQ$ 是否经过 $x$ 轴上的定点？若过定点，求出该点坐标；若不过定点，说明理由.
\end{example}

\begin{example} % 直线，三角形面积
    % 练习题目
（2023海淀一模）已知椭圆 $E:\frac{x^2}{a^2}+\frac{y^2}{b^2}=1\ (a>b>0)$ 的左、右顶点分别为 $A_1,A_2$，上、下顶点分别为 $B_1,B_2$，$\lvert B_1B_2\rvert=2$，四边形 $A_1B_1A_2B_2$ 的周长为 $4\sqrt{6}$.\\
（1）求椭圆 $E$ 的方程；\\
（2）设斜率为 $k$ 的直线 $l$ 与 $x$ 轴交于点 $P$，与椭圆 $E$ 交于不同的两点 $M,N$，点 $M$ 关于 $y$ 轴的对称点为 $M'$，直线 $M'N$ 与 $y$ 轴交于点 $Q$. 若 $\triangle OPQ$ 的面积为 $2$，求 $k$ 的值.
\end{example}

\begin{example} % 动点
（2022丰台一模）已知椭圆 $C:\frac{x^2}{a^2}+\frac{y^2}{b^2}=1\ (a>b>0)$ 的左、右顶点分别为 $A,B$，且 $\lvert AB\rvert=4$，离心率为 $\frac{\sqrt{3}}{2}$.\\
（1）求椭圆 $C$ 的方程；\\
（2）设 $P$ 是椭圆 $C$ 上不同于 $A,B$ 的一点，直线 $PA,PB$ 与直线 $x=4$ 分别交于点 $M,N$. 若 $\lvert MN\rvert\leq4$，求点 $P$ 横坐标的取值范围.
\end{example}

\begin{example} % 直线引入，三点共线，定值定点（不需要什么技巧）
    % 练习题
（2022海淀一模）已知椭圆 $C:\frac{x^2}{a^2}+\frac{y^2}{b^2}=1\ (a>b>0)$ 的下顶点 $A$ 和右顶点 $B$ 都在直线 $l_1:y=\frac{1}{2}(x-2)$ 上.\\
（1）求椭圆方程及其离心率；\\
（2）不经过点 $B$ 的直线 $l_2:y=kx+m$ 交椭圆 $C$ 于两点 $P,Q$，过点 $P$ 作 $x$ 轴的垂线交 $l_1$ 于点 $D$，点 $P$ 关于点 $D$ 的对称点为 $E$. 若 $E,B,Q$ 三点共线，求证：直线 $l_2$ 经过定点.
\end{example}

\begin{example} % 直线引入，定值定点与非对称韦达
    % 非对称与定值例题
（2022朝阳一模）已知椭圆 $C:\frac{x^2}{a^2}+\frac{y^2}{b^2}=1\ (a>b>0)$ 的一个焦点为 $F(1,0)$，且过点 $\left(1,\frac{3}{2}\right)$.\\
（1）求椭圆 $C$ 的方程和离心率；\\
（2）过点 $P(4,0)$ 且与 $x$ 轴不重合的直线 $l$ 与椭圆 $C$ 交于 $A,B$ 两点，与直线 $x=1$ 交于点 $Q$，点 $M$ 满足 $MP\perp x$ 轴，$MB\parallel x$ 轴，试求直线 $MA$ 的斜率与直线 $MQ$ 的斜率的比值.
\end{example}

\begin{example} % 直线，角度转长度
（2022西城一模）已知椭圆 $C:\frac{x^2}{a^2}+\frac{y^2}{b^2}=1\ (a>b>0)$ 的离心率为 $\frac{\sqrt{3}}{2}$，以椭圆的四个顶点为顶点的四边形周长为 $4\sqrt{5}$.\\
（1）求椭圆的方程；\\
（2）直线 $y=kx+m\ (km\neq0)$ 与椭圆 $C$ 交于 $A,B$ 两点，与 $y$ 轴交于点 $P$，线段 $AB$ 的垂直平分线与 $AB$ 交于点 $M$，与 $y$ 轴交于点 $N$，$O$ 为坐标原点. 如果 $\angle MOP=2\angle MNP$，求 $k$ 的值.
\end{example}

\begin{example} % 直线，弦长公式
    % 练习题
（2022东城一模）已知椭圆 $C:\frac{x^2}{a^2}+\frac{y^2}{b^2}=1\ (a>b>0)$ 的离心率为 $\frac{\sqrt{3}}{2}$，焦距为 $2\sqrt{3}$.\\
（1）求椭圆 $C$ 的方程；\\
（2）过点 $P(4,0)$ 作斜率为 $k$ 的直线 $l$ 与椭圆 $C$ 交于 $A,B$ 两点. 是否存在常数 $t$，使得直线 $x=t$ 与直线 $l$ 的交点 $Q$ 在 $A,B$ 之间，且总有 $\frac{\lvert PA\rvert}{\lvert PB\rvert}=\frac{\lvert QA\rvert}{\lvert QB\rvert}$？若存在，求出 $t$ 的值；若不存在，说明理由.
\end{example}

\begin{example} % 动点，定值定点
    % 定点例题
（2021海淀一模）已知椭圆 $M:\frac{x^2}{a^2}+\frac{y^2}{b^2}=1\ (a>b>0)$ 过 $A(-2,0),B(0,1)$ 两点.\\
（1）求椭圆 $M$ 的离心率；\\
（2）设椭圆 $M$ 的右顶点为 $C$，点 $P$ 在椭圆 $M$ 上（$P$ 不与椭圆 $M$ 的顶点重合），直线 $AB$ 与直线 $CP$ 交于点 $Q$，直线 $BP$ 交 $x$ 轴于点 $S$，求证：直线 $SQ$ 过定点.
\end{example}

\begin{example} % 直线，定值
（2021东城一模）已知椭圆 $C:\frac{x^2}{a^2}+\frac{y^2}{b^2}=1\ (a>b>0)$ 过点 $D(-2,0)$，且焦距为 $2\sqrt{3}$.\\
（1）求椭圆 $C$ 的方程；\\
（2）过点 $A(-4,0)$ 的直线 $l$（不与 $x$ 轴重合）与椭圆 $C$ 交于 $P,Q$ 两点，点 $T$ 与点 $Q$ 关于 $x$ 轴对称，直线 $TP$ 与 $x$ 轴交于点 $H$. 是否存在常数 $\lambda$，使得 $\lvert AD\rvert\cdot\lvert DH\rvert=\lambda(\lvert AD\rvert-\lvert DH\rvert)$ 成立？若存在，求出 $\lambda$ 的值；若不存在，说明理由.
\end{example}

\begin{example} % 直线，定点，非对称韦达
    % 定点例题
（2021西城一模）已知椭圆 $C:\frac{x^2}{a^2}+\frac{y^2}{3}=1\ (a>0)$ 的焦点在 $x$ 轴上，且经过点 $E\left(1,\frac{3}{2}\right)$，左顶点为 $D$，右焦点为 $F$.\\
（1）求椭圆 $C$ 的离心率和 $\triangle DEF$ 的面积；\\
（2）已知直线 $y=kx+1$ 与椭圆 $C$ 交于 $A,B$ 两点，过点 $B$ 作直线 $y=t\ (t>\sqrt{3})$ 的垂线，垂足为 $G$. 判断是否存在常数 $t$，使得直线 $AG$ 经过 $y$ 轴上的定点？若存在，求 $t$ 的值；若不存在，请说明理由.
\end{example}

\begin{example} % 动点，以线段为直径的圆过定点，定点，特殊法先得答案
    % 练习题
（2021朝阳一模）已知椭圆 $C$ 的短轴的两个端点分别为 $A(0,1),B(0,-1)$，离心率为 $\frac{\sqrt{6}}{3}$.\\
（1）求椭圆 $C$ 的方程及焦点的坐标；\\
（2）若点 $M$ 为椭圆 $C$ 上异于 $A,B$ 的任意一点，过原点且与直线 $MA$ 平行的直线与直线 $y=3$ 交于点 $P$，直线 $MB$ 与直线 $y=3$ 交于点 $Q$. 试判断以线段 $PQ$ 为直径的圆是否过定点？若过定点，求出定点的坐标；若不过定点，请说明理由.
\end{example}

\begin{example} % 动点，三点共线
（2021丰台一模）已知椭圆 $C:\frac{x^2}{a^2}+\frac{y^2}{b^2}=1\ (a>b>0)$ 长轴的两个端点分别为 $A(-2,0),B(2,0)$，离心率为 $\frac{\sqrt{3}}{2}$.\\
（1）求椭圆 $C$ 的方程；\\
（2）$P$ 为椭圆 $C$ 上异于 $A,B$ 的动点，直线 $AP,PB$ 分别交直线 $x=-6$ 于 $M,N$ 两点，连接 $NA$ 并延长交椭圆 $C$ 于点 $Q$.\\
（i）求证：直线 $AP,AN$ 的斜率之积为定值；\\
（ii）判断 $M,B,Q$ 三点是否共线，并说明理由.
\end{example}



